\documentclass[10pt,conference,compsocconf]{IEEEtran}
\usepackage{amsmath}
\usepackage{url}
\usepackage[hidelinks]{hyperref}

\begin{document}

\title{Predicting Coronary Heart Disease from BRFSS 2015 with\\
Plain SGD and L2-Regularised Logistic Regression}

\author{
  \IEEEauthorblockN{Youssef Ahmed}
  \IEEEauthorblockA{\textit{Westlake University}\\
  \textit{Machine Learning Course --- Project 1}}
}

\maketitle

\begin{abstract}
We predict reported coronary heart disease, myocardial infarction, or angina
from the 2015 BRFSS (328{,}135 labeled; 8.83\% positive; 44.8\% missing). We
use plain single-sample SGD logistic regression with an all-weight L2 penalty.
On the seed-42 training-side partition (also used earlier), $\lambda=10^{-5}$
had the best calibration F1 and a 27\% smaller norm than $\lambda=0$. Full-data
5-fold CV gives $0.4303\pm0.0060$ F1, 0.8622 AUROC and 0.4188 PR-AUC;
unregularized encoded gives $0.4315\pm0.0053$ F1. The frozen threshold is
0.2148; final predictions are 17{,}310 positive.
\end{abstract}

\section{Introduction and Data}
Coronary heart disease is a leading cause of death, and survey data can help
identify respondents with a reported history of heart disease. We use the 2015
BRFSS~\cite{brfss15}; the target
\texttt{\_MICHD} is positive when a respondent reports that a health
professional told them they had coronary heart disease, a myocardial
infarction, or angina. The training set has 328{,}135 labeled rows (28{,}975
positive, 8.83\%), 321 raw features, substantial missingness and coded
categorical variables. The test set has 109{,}379 rows.

\section{Preprocessing and Feature Representation}
For every fitting portion, medians, means, standard deviations and category
sets are learned from that portion only and applied unchanged to held-out or
test data. We remove 14 identifiers and survey-design fields (interview date
components, record number, disposition code, sampling unit, stratum and survey
weights), remove constant columns, median-impute and standardise numeric
features, and one-hot encode features with at most 15 distinct finite training
values. We use $K-1$ dummies, the most frequent value as reference, and an
explicit missing category. Designs have 1{,}129--1{,}130 columns on the
262{,}508-row training portions and 1{,}120 on all data.

A header-level check confirmed that \texttt{\_MICHD} is the separate target and
that \texttt{CVDINFR4} and \texttt{CVDCRHD4}, direct source variables for its
definition, are absent from the predictors; this does not prove the absence of
every derived relationship. We retained \texttt{CVDSTRK3}, which records stroke
history rather than the target condition.

\section{Six Required Course Algorithms}
\texttt{implementations.py} contains the six required functions:
\texttt{mean\_squared\_error\_gd}, \texttt{mean\_squared\_error\_sgd},
\texttt{least\_squares}, \texttt{ridge\_regression},
\texttt{logistic\_regression} and \texttt{reg\_logistic\_regression}. Their
signatures and behavior are unchanged, and they pass all 25 public tests. The
final model and the experiments use the shared plain-SGD routine in
\texttt{pipeline.py}; \texttt{run.py} is the final prediction entry point.

\section{Plain Stochastic Gradient Descent}
For labels $y_i\in\{0,1\}$ and logits $z_i=x_i^{\top}w$, the mean logistic loss
is
\begin{equation}
L(w)=\frac1N\sum_{i=1}^{N}\left[\log(1+e^{z_i})-y_i z_i\right].
\end{equation}
For one example, the update is
\begin{equation}
w\leftarrow w-\eta_t\left(\sigma(x_i^{\top}w)-y_i\right)x_i,
\end{equation}
where $\sigma(z)=1/(1+e^{-z})$ and $t$ is the epoch index. We use the inverse
schedule $\eta_t=0.01/(1+0.1(t-1))$ for 90 epochs, reshuffling each epoch with
seed 2027. There are no moment estimates or per-parameter scaling. We use a
stable sigmoid and log loss and compute the full loss only at checkpoints.

We verified the analytic gradient against central finite differences for the
unregularized and L2 objectives, the average of per-sample gradients against
the full-data gradient, a manual update, schedule values, seed-reproducible
shuffling, finite losses, and the absence of adaptive state. The loss decreased
on a controlled synthetic problem. Because SGD is noisy, we use a fixed epoch
cap rather than validation early stopping.

\section{L2 Regularisation}
The penalized objective is $J(w)=L(w)+\lambda\sum_{j=0}^{d}w_j^2$, so
\begin{equation}
w\leftarrow w-\eta_t\left[
\left(\sigma(x_i^{\top}w)-y_i\right)x_i+2\lambda w\right].
\end{equation}
The penalty includes the intercept, matching the course convention. We studied
$\lambda\in\{0,10^{-5},10^{-4},10^{-3}\}$ on the seed-42 stratified 80/20
partition of all labeled data: 262{,}508 fit and 65{,}627 calibration rows.
This partition was also used during earlier development. The fitting rows train
each $\lambda$; the calibration rows select the threshold and supply the F1
values, so this study describes model selection and regularization behavior,
not independent generalization. It is not the five-fold CV; the same SGD
configuration and epochs were used for every $\lambda$.

\begin{table}[t]
\centering
\small
\caption{L2 study on the seed-42 training-side 80/20 partition (262{,}508 fit,
65{,}627 calibration), also used earlier; F1 uses a calibration-optimized
threshold, not independent validation. $\Delta$ is calibration minus training loss.}
\label{tab:reg}
\setlength{\tabcolsep}{3pt}
\begin{tabular}{lcccc}
\hline
$\lambda$ & F1 & AUROC & $\|w\|_2$ & $\Delta$\\
\hline
0 & 0.4368 & 0.8622 & 10.537 & $+0.00219$\\
$10^{-5}$ & \textbf{0.4374} & \textbf{0.8625} & 7.663 & $+0.00182$\\
$10^{-4}$ & 0.4360 & 0.8623 & 4.180 & $+0.00109$\\
$10^{-3}$ & 0.4284 & 0.8581 & 2.060 & $+0.00062$\\
\hline
\end{tabular}
\end{table}

The weakly regularized model achieved a slightly lower calibration logistic
loss (0.22021 versus 0.22037) and a 27\% smaller norm. $\lambda=10^{-4}$ shrinks
the norm by 60\% at a cost of 0.0014 F1; $\lambda=10^{-3}$ underfits (F1 0.4284,
AUROC 0.8581). The training--calibration gap narrows with stronger penalties.
A smaller norm alone is not evidence of better generalization, and the CV shows
no uniform F1 advantage for the regularized model; $\lambda=10^{-5}$ is retained
because the calibration study supports it and coefficients are controlled.

\section{Full-Data Stratified Five-Fold Cross-Validation}
\label{sec:cv}
We run one stratified five-fold CV on all 328{,}135 labeled examples
(28{,}975 positives, 299{,}160 negatives). Each fold holds out exactly 65{,}627
rows: 5{,}795 positives and 59{,}832 negatives. Every row appears in exactly one
evaluation fold and four training sets; the fold seed is 2028 and assignments
are saved.

Each 262{,}508-row outer training portion is split, stratified, into 210{,}006
fitting and 52{,}502 inner calibration rows. Preprocessing is fitted only on the
fitting rows. The threshold is chosen on inner calibration data to maximize
positive-class F1 and frozen; the model and preprocessing are refitted on the
complete outer training portion, and the threshold is applied once to the
held-out fold. No held-out label influences its own preprocessing, threshold,
or fit. All models use the same folds.

Hyperparameters, $\lambda$, and the final threshold were frozen before this CV;
all labeled rows had participated in earlier training-side development, so the
CV is a retrospective evaluation of a frozen procedure, not an untouched
independent test. The folds are correlated and their SD is descriptive.

\begin{table}[t]
\centering
\small
\caption{Full-data stratified five-fold CV on all 328{,}135 labeled examples.
F1 is mean $\pm$ sample SD over five held-out folds; AUROC and PR-AUC are fold
means. Thresholds are selected on inner data only.}
\label{tab:cv}
\setlength{\tabcolsep}{3pt}
\begin{tabular}{lccc}
\hline
Model & F1 (mean $\pm$ SD) & AUROC & PR-AUC\\
\hline
Majority ($-$1) & 0.0000 $\pm$ 0.0000 & 0.5000 & 0.0883\\
Ridge (encoded) & 0.4272 $\pm$ 0.0027 & 0.8583 & 0.4161\\
Logistic (numeric) & 0.4022 $\pm$ 0.0036 & 0.8411 & 0.3502\\
Logistic (encoded) & 0.4315 $\pm$ 0.0053 & 0.8619 & 0.4188\\
\textbf{L2 logistic (encoded)} & \textbf{0.4303 $\pm$ 0.0060} & \textbf{0.8622} & \textbf{0.4188}\\
\hline
\end{tabular}
\end{table}

The unregularized encoded model has the largest mean F1 (0.4315); the selected
L2 model is 0.0012 lower. The paired difference, F1(L2) $-$ F1(unregularized),
is $-0.0012\pm0.0029$ and changes sign across folds. L2 has slightly higher
AUROC and the same PR-AUC. We do not claim a significant or uniform advantage;
L2 is retained because the separate calibration study selected it, it controls
coefficient magnitude, and full-data F1 is comparable.

The encoded model's paired mean F1 advantage over numeric codes is
$+0.0293\pm0.0043$ (encoded minus numeric); over ridge it is
$+0.0043\pm0.0051$, and ridge has lower mean AUROC. These differences are
descriptive, not significance tests.

\section{Threshold Selection and Final Model}
The final threshold is $0.21475973018819802$, selected to maximize positive-class
F1 on an inner split of the 262{,}508-row seed-42 training portion (210{,}006
fit, 52{,}502 calibration, inner seed 123; inner F1 0.4336) and frozen before
the all-data fit. Each CV
fold selects its own threshold (0.1538--0.2427). The regularization study
instead fits on all 262{,}508 rows and optimizes its own threshold (about
0.1706, F1 about 0.4374); applying the frozen threshold to those scores gives
F1 about 0.4255. These are different, non-independent procedures.

The final L2 model is trained on all 328{,}135 labeled rows with the frozen
configuration and applied to the 109{,}379 test rows. It produces 17{,}310
positive predictions (15.83\%). Test labels are unavailable, so no test
performance is reported. The final training loss is 0.2189.

\section{Discussion and Limitations}
The model is linear and cannot capture non-linear interactions; no resampling
or class weighting was used. The full-data CV is retrospective because all
labeled rows informed earlier preprocessing, hyperparameter, and threshold
choices, and its folds are correlated. The 90-epoch cap and schedule were
frozen before CV; we do not claim convergence. The L2 experiment measures
coefficient shrinkage and calibration-partition loss, not dedicated probability
calibration. Associations are not causal, the threshold is a point estimate,
and test labels are unavailable.

\section{Conclusion}
A NumPy-only pipeline combines one-hot encoding, verified preprocessing, plain
single-sample SGD, an all-weight L2 penalty, and leakage-free fold-level
threshold selection. Full-data five-fold CV gives $0.4303\pm0.0060$ F1, 0.8622
AUROC and 0.4188 PR-AUC; the unregularized encoded model is 0.0012 higher, a
small and mixed difference. The final model produces 17{,}310 positive test
predictions.

\begin{thebibliography}{1}
\bibitem{brfss15}
Centers for Disease Control and Prevention.
\textit{2015 BRFSS Survey Data and Documentation}.
\url{https://www.cdc.gov/brfss/annual_data/annual_2015.html}, 2016.
\end{thebibliography}

\end{document}
